@ID: ON21D1P1U
@BIDANG: Kombinatorika
Diberikan bilangan bulat positif $n$ dan $k$ sehingga $2\le k\le n$. Buktikan bahwa

$$
\sum_{k=2}^{n}\binom{k}{2}\binom{n+2}{k}\binom{n-2}{k-2}=\binom{n+2}{2}\binom{2n-2}{n-2}
$$

 208].

@ID: ON21D1P2U
@BIDANG: Analisis Kompleks
Diketahui fungsi $f$ analitik di dalam dan pada kurva tertutup sederhana $\gamma$ dan $z_{0}$ sebuah titik pada bidang kompleks yang tidak berada pada $\gamma$. Buktikan untuk setiap bilangan asli $m$ dan $n$ berlaku

$$
\oint_{\gamma}\frac{f^{(m)}(z)}{(z-z_{0})^{n}}dz=\frac{(m+n-1)!}{(n-1)!}\oint_{\gamma}\frac{f(z)}{(z-z_{0})^{m+n}}dz.
$$

 211].

@ID: ON21D1P3U
@BIDANG: Struktur Aljabar
Misalkan $k$ dan $n$ bilangan bulat positif dan $G$ adalah grup dengan $n$ unsur. Buktikan bahwa kedua pernyataan berikut ekivalen.

(a) Bilangan $n$ dan $k$ relatif prima.

(b) Untuk setiap subgrup $H$ dari $G$ berlaku $\{x\in G|x^{k}\in H\}\subseteq H$.

@ID: ON21D1P4U
@BIDANG: Aljabar Linear
Diberikan matriks $A$ berukuran $n\times n$ dengan komponen-komponen real. Didefinisikan pemetaan linear $T:\mathbb{R}^{n}\rightarrow\mathbb{R}^{n}$ dan $S:\mathbb{R}^{n}\rightarrow\mathbb{R}^{n}$ dengan $T(x)=Ax$ dan $S(x)=A^{T}x$ untuk setiap $x\in\mathbb{R}^{n}$.

(a) Buktikan bahwa $\text{Ker}(S\circ T)=\text{Ker}(T)$.

(b) Jelaskan apakah selalu berlaku $\text{Im}(S\circ T)=\text{Im}(S)$! 
Catatan: $\text{Ker}(T)$ menyatakan kernel dari $T$ dan $\text{Im}(T)$ menyatakan peta dari $T$.

@ID: ON21D1P5U
@BIDANG: Analisis Real
Diberikan barisan fungsi kontinu $(f_{n})$ dengan $f_{n}:[0,1]\rightarrow[0,\infty)$, $\forall n\in\mathbb{N}$. Jika dipenuhi kondisi berikut:
(i) $f_{1}(x)\ge f_{2}(x)\ge f_{3}(x)\ge\dots\forall x\in[0,1]$;
(ii) $f(x)=\lim_{n\rightarrow\infty}f_{n}(x)$ dan $M=\sup_{0\le x\le1}f(x)$ 226],
maka:

(a) buktikan terdapat $t\in[0,1]$ sedemikian sehingga $f(t)=M$.

(b) berikan contoh bahwa 5(a) belum tentu berlaku, apabila kondisi (i) diganti menjadi: terdapat $N_{x}\in\mathbb{N}$ sehingga $\forall n\ge N_{x}$, $f_{n}(x)\ge f_{n+1}(x)$. Jelaskan!.

@ID: ON21D2P1U
@BIDANG: Kombinatorika
Tentukan banyaknya bilangan terurut quintuples $(a_{1},a_{2},a_{3},a_{4},a_{5})$ dari bilangan bulat ganjil positif yang memenuhi $a_{1}+a_{2}+a_{3}+a_{4}+a_{5}=101$.

@ID: ON21D2P2U
@BIDANG: Analisis Real
Diberikan himpunan terbatas $A\subset\mathbb{R}$. Jika fungsi $f:A\rightarrow\mathbb{R}$ kontinu seragam,

(a) selidiki apakah terdapat fungsi kontinu $g:\overline{A}\rightarrow\mathbb{R}$ dengan $g(x)=f(x)$, untuk setiap $x\in A$.

(Catatan: $\overline{A}$ menyatakan klosur himpunan $A$) .

(b) buktikan bahwa $f$ terbatas.

@ID: ON21D2P3U
@BIDANG: Aljabar Linear
Untuk sebarang bilangan asli $n$, misalkan $A_{n}=(a_{ij})$ matriks berukuran $n\times n$ dengan $a_{ij}=|i-j|$ untuk setiap $1\le i, j\le n$. Tentukan, dengan bukti, nilai dari $\det(A_{n})$ untuk setiap bilangan asli $n$.

@ID: ON21D2P4U
@BIDANG: Analisis Kompleks
Pasangan terurut bilangan real $(a, b)$ dikatakan milenial jika memenuhi persamaan $(a+bi)^{2021}=a^{2}-b^{2}-2abi$ dengan $i=\sqrt{-1}$.

(a) Jika $(a, b)$ adalah pasangan milenial, tunjukan bahwa $a=b=0$ atau $a^{2}+b^{2}=1$.

(b) Hitunglah banyaknya semua pasangan milenial.

@ID: ON21D2P5U
@BIDANG: Struktur Aljabar
Misalkan $(A,+)$ adalah grup siklis dengan kardinalitas $n$. Ditetapkan $a$ sebuah pembangun dari $(A, +)$. Untuk setiap $m$ bilangan asli, didefinisikan ring $A_{m}$ dengan $(A_{m},+)=(A,+)$ dan perkalian di $A_{m}$ didefinisikan melalui

$$
a\cdot a=\underbrace{a+a+\dots+a}_{m \text{ kali}} = ma
$$

Buktikan bahwa $A_{r}$ isomorfik dengan $A_{s}$ sebagai ring jika dan hanya jika $\text{FPB}(r,n)=\text{FPB}(s,n)$.